% ================================================================
%                          译者说明
% ================================================================
\chapter*{译者说明}
\addcontentsline{toc}{chapter}{译者说明}
\markboth{译者说明}{译者说明}

本译本依据 Tobias Holck Colding 与 William P. Minicozzi II 合著的
\emph{A Course in Minimal Surfaces}(American Mathematical Society,
Graduate Studies in Mathematics 第 121 卷, 2011 年初版)译出并重排. 原书共八章:
第 1 章从面积泛函的 Euler--Lagrange 方程出发建立极小曲面的基本理论;
第 2 章讨论 Simons 不等式及其推论, 包括各种曲率估计与 Bernstein 型定理;
第 3 章引入变分流形(广义极小曲面)与 Sobolev 不等式, 讨论弱收敛与紧性;
第 4 章给出 Plateau 问题与 Dirichlet 问题的存在性理论;
第 5 章用极小极大方法构造极小球面;
第 6 章研究 Plateau 问题解的正则性与嵌入性;
第 7 章把上述理论推广到三维流形中的极小曲面;
第 8 章介绍嵌入极小曲面结构理论的最新结果. 全书之后附有习题、参考文献与索引.

\textbf{翻译原则.} 本译本遵循数译(mathtranslations.org)的《数学译本制作规范》:
忠于原文, 数学正确优先于语言流畅, 术语全书统一, 保留原书的数学结构(定义、
定理、引理、推论、命题、注、证明的层次与编号), 并尽可能实现可点击的交叉引用.
正文中文用宋体, 首次出现的专门术语以楷体标注并附英文原词; 定理类正文用仿宋,
例子与问题用宋体; 英文与数学符号用衬线 Computer Modern 字体. 除句子末尾使用
英文句点外, 其余标点均用中文标点.

\textbf{编号说明.} 原书的带编号小标题(如``1.1. Graphs and the minimal surface
equation.'')与定理类对象共用同一个按章推进的计数器, 因此同一个序号的类型可能
不同(``1.1'' 是小标题, ``1.5'' 是引理). 本译本保留这一版式: 定理、引理、推论、
命题、定义、注、例与带编号小标题共享计数器, 显示为``章号.序号''. 本书自身的公式、定理、图与节的引用均经交叉引用生成, 不手写编号;
引自其他文献的定位符(如``\cite{GiTr} 的定理 8.8'')是原书对别的书的编号,
保留原书的手写形式. 原书个别错引(例如某处写作``定理 5.5''而 5.5 实为引理)
一律照录, 不做改正, 也不加链接.

\textbf{图形说明.} 原书为扫描版, 书中插图无法从矢量源重新生成, 故按原书页面
裁切后嵌入, 分辨率与原书扫描精度一致, 并逐图核对裁切范围.

\textbf{索引说明.} 原书末尾的 Index 以主题词加页码的形式编排. 由于本译本的
分页与原书不同, 逐条照搬页码没有意义, 故以模板提供的术语索引重建: 每个术语
在其\textbf{首次出现处}登记, 索引中的页码指向本译本中的对应位置, 并按术语在
正文中首次出现的顺序排列.

\textbf{模型使用说明.} 本译本由数译团队使用大语言模型辅助翻译, 数学公式与
符号依照原书排版逐式还原; 原书扫描中的个别模糊字符已对照页面图像人工核定.

\vspace{0.8cm}
\hfill \Translator

\hfill \TranslationDate

\cleardoublepage
